% Fail the warm-cache CI run if any CJK face needs raw OpenType parsing.
\RequirePackage{luaotfload}
\directlua{
  if os.getenv("REQUIRE_PREWARMED_FONTS") == "1" then
    local readers = fonts.handlers.otf.readers
    local loadfont = readers.loadfont
    readers.loadfont = function(filename, ...)
      if string.find(filename, "HaranoAji", 1, true) then
        error("PREWARMED_FONT_CACHE_MISS")
      end
      return loadfont(filename, ...)
    end
  end
}
\documentclass[a4paper,11pt]{ltjsarticle}
\usepackage{luatexja-fontspec}
\setmainjfont{Harano Aji Mincho}
\setsansjfont{Harano Aji Gothic}
\usepackage{amsmath,booktabs,tikz}
\begin{document}
\section{フォントキャッシュの検証}
明朝体、日本語の本文です。\textbf{太字の日本語。}
\textsf{ゴシック体の日本語。}\textsf{\textbf{ゴシック体の太字。}}
\begin{equation} \int_0^1 x^2\,dx = \frac{1}{3} \end{equation}
\begin{align} a^2+b^2&=c^2 \\ e^{i\pi}+1&=0 \end{align}
\begin{tikzpicture}
\draw[->] (0,0)--(3,0);
\draw[->] (0,0)--(0,2);
\draw (0,0)--(2,1);
\node at (1,1.5) {日本語};
\end{tikzpicture}
\newpage
\section{二ページ目}
本文と数式の再現性を確認します。\(\sum_{n=1}^{10}n=55\)
\end{document}
